\begin{abstract}
We introduce Lotka-Volterra Competitive Serving (LVCS), a novel dynamic model ensembling paradigm that rejects traditional linear state-space formulations in favor of non-linear biological ecosystem modeling. Drawing inspiration from population ecology, LVCS models the activation-space trajectories of task-specific experts as population densities of competing species governed by a discrete-time Lotka-Volterra Ricker recurrence. We introduce learned diagonal carrying capacities to stabilize representation flow under high activation noise, alongside Adaptive Niche Plasticity to dynamically suspend inter-species competition during sudden stream transitions, thereby eliminating representational lag. Furthermore, we propose a systems-first static coordinate approximation that extracts routing coordinates once at an early layer, reducing serving latency by over 51\% compared to a fully dynamic model while maintaining virtually identical accuracy. Our comprehensive evaluation demonstrates that LVCS achieves stable and responsive routing, outperforming the state-of-the-art linear stateful baseline (PAC-Kinetics) by up to $+1.28\%$ absolute accuracy on overlapping task manifolds while completely resolving representational lag under heterogeneous sequential serving streams. This work establishes a formal, biologically-grounded bridge between mathematical ecology and systems-efficient dynamic model ensembling.
\end{abstract}
